\section{The four-pole aggregate and every cyclic four-edge cut}
\label{sec:all-four-cuts}

For a rank-$h$ four-pole, use normalized path counts $a_{ij}$ and normalized
forest-pairing counts $u_\alpha$ as in Section~\ref{sec:four-cut-large-shores}.
Put
\[
 S=\sum_{i<j}a_{ij},\qquad F=\sum_\alpha u_\alpha.
\]
Each path count and $F$ lies in $[0,1]$.

\begin{lemma}[The added-edge constraint]\label{lem:four-pole-added-edge}
For each $0<\delta,\epsilon<1$, there is an absolute rank threshold
$H_{\delta,\epsilon}$ such that the following holds.
Assume the four-pole rank satisfies $h\ge H_{\delta,\epsilon}$.
Let $\alpha=ij|kl$ be one of the three opposite pairings.
If $a_{kl}\ge\delta$, then
\[
 a_{ij}+F-u_\alpha\le1+\epsilon.
\]
If both $a_{ij},a_{kl}\ge\delta$, then
\[
 \max(a_{ij},a_{kl})+F-u_\alpha\le1+\epsilon.
\]
\end{lemma}

\begin{proof}
Add an edge $kl$ to the four-pole.
This gives a normalized two-pole $N$ with terminals $i,j$ and rank $h+1$.
Its terminal-path and cycle counts satisfy
\[
 P_{ij}(N)=2^h(a_{ij}+F-u_\alpha),\qquad c(N)\ge a_{kl}2^h.
\]
Paths using the new edge correspond exactly to the two cross forest pairings.
Cycles using it correspond to the original $k$--$l$ paths.
If the asserted inequality failed, the normalized path density of $N$
would exceed $1/2+\epsilon/2$.
Theorem~\ref{thm:high-path-cycle-tradeoff} would imply
\[
 \delta2^h\le c(N)\le K_{\epsilon/2}2^{(h+1)/2},
\]
which is impossible for sufficiently large $h$.
Reverse the roles of the opposite pairs to obtain the second assertion.
\end{proof}

\begin{theorem}[The four-pole aggregate]\label{thm:four-pole-aggregate}
There is an absolute constant $A_4$ such that every connected loopless
cubic four-pole of rank $h$ satisfies
\[
 \sum_{i<j}P_{ij}+\sum_\alpha f_\alpha
 \le4\,2^h+A_4(h+3)^4 2^{h/2}.
\]
\end{theorem}

\begin{proof}
Suppress each unused fourth port and apply
Theorem~\ref{thm:universal-tripole-sum} to the other three.
Summing the four resulting inequalities counts each pair twice.
Consequently, for an absolute constant $C$,
\begin{equation}\label{eq:four-pole-total-paths}
 S\le4+C2^{-h/2}.
\end{equation}
If suppression creates a loop, every path between the remaining terminals
avoids its cycle and has density at most one half.
The relevant three-terminal sum bound therefore remains valid in that case.

If some cross sum is at least $2.01$, the cross-sum forest bound gives
$F\le C'(h+3)^4 2^{-h/2}$, proving the result.
Assume that every cross sum is below $2.01$.
The three cross sums add to $2S$, so $S<3.015$.
There is nothing to prove if $S+F\le4$.
Otherwise $S>3$ and $F>0.985$.

For each opposite pairing $\alpha$, let $x_\alpha$ be the sum of its two
path densities. Its complementary cross sum is $S-x_\alpha$.
Thus
\begin{equation}\label{eq:four-pole-opposite-positive}
 x_\alpha>0.99\qquad\hbox{for every }\alpha.
\end{equation}
Suppose first that no path pair is full.
The terminal-path gap bounds each density by $15/16$.
Equation~\eqref{eq:four-pole-opposite-positive} then makes every density
greater than $0.99-15/16=0.0525$.
Apply Lemma~\ref{lem:four-pole-added-edge} with $\delta=\epsilon=0.01$.
Writing $M_\alpha$ for the larger density in opposite pair $\alpha$, summation gives
\[
 \sum_\alpha M_\alpha+2F\le3.03.
\]
Since $S\le2\sum_\alpha M_\alpha$, we get $S+4F\le6.06$.
This contradicts $S>3$ and $F>0.985$.

Now suppose that $a_{12}=1$.
Lemma~\ref{lem:one-full-tripole-small-adjacent}, applied after suppressing
the unused fourth port, gives one density at most $3\,2^{-h/4}$ in each set
\[
 \{a_{13},a_{23}\},\qquad \{a_{14},a_{24}\}.
\]
For sufficiently large $h$, these two densities are at most $0.01$.
Their opposite densities exceed $0.98$ by
\eqref{eq:four-pole-opposite-positive}, and hence are full by the $15/16$ gap.
The two small densities cannot themselves be opposite, since their sum
would violate~\eqref{eq:four-pole-opposite-positive}.
Therefore the three full pairs form a star.
Theorem~\ref{thm:three-pole-two-high} bounds each of its three leaf-pair
densities by $C''2^{-h/2}$.
It follows that $S+F\le4+3C''2^{-h/2}$.

All thresholds used above are absolute.
For the finitely many smaller ranks, the trivial bound $S+F\le7$ is
absorbed by increasing $A_4$.
\end{proof}

\begin{lemma}[Low path states]\label{lem:four-pole-low-state}
For sufficiently large $h$, suppose all six normalized path densities
of a four-pole are at most $\tau=0.51$.
Then
\[
 S+F\le3.82.
\]
For any selected pairing $\gamma$, one also has
\[
 S+\frac12F+\frac12u_\gamma\le3.7.
\]
\end{lemma}

\begin{proof}
If $S\le2.75$, the first bound follows from $F\le1$.
Otherwise every path density exceeds $2.75-5\tau=0.20$.
Apply the added-edge constraint with $\delta=\epsilon=0.01$ to all three pairings.
It gives $S/2+2F\le3.03$.
Since $S\le6\tau=3.06$,
\[
 S+F\le1.515+\frac34S\le3.81<3.82.
\]

For the second bound, the case $S\le2.7$ follows from $u_\gamma\le F\le1$.
Otherwise every density exceeds $2.7-5\tau=0.15$.
Let $\alpha,\beta$ be the other pairings.
The two corresponding added-edge constraints give
\[
 M_\alpha+M_\beta+F+u_\gamma\le2.02.
\]
Consequently,
\[
 S+\frac12(F+u_\gamma)
 \le S-\frac12(M_\alpha+M_\beta)+1.01.
\]
The two coordinates in each of $\alpha,\beta$ contribute at most
$3\tau/2$ after this subtraction.
The coordinates in $\gamma$ contribute at most $2\tau$.
The last display is therefore at most $5\tau+1.01=3.56<3.7$.
\end{proof}

\begin{theorem}[Every cyclic cut of size at most four]
\label{thm:all-four-cut-half-exponent}
There is an absolute constant $K_4$ such that every connected loopless
cubic multigraph with a cyclic edge cut of size at most four satisfies
\[
 c(G)\le2^{r-1}+K_4r^8 2^{r/2},\qquad r=\beta(G).
\]
A cut is cyclic when both shores contain a cycle.
\end{theorem}

\begin{proof}
Cyclic cuts of size at most three are covered by
Theorem~\ref{thm:all-three-cut-half-exponent}.
Assume none exists, and consider a cyclic cut of size four.
The graph has no bridge or two-edge cut, since those are cyclic in a cubic graph.
Thus each component after deleting the four cut edges has boundary at least three.
There are exactly two components: three would require total boundary at least nine,
whereas the four cut edges supply only eight.
Write these connected cyclic shores as $X,Y$.

Their attachment vertices are distinct.
Indeed, a vertex with two cut edges has one internal edge and is a leaf in its shore.
Move this vertex to the other shore.
The original shore remains connected and cyclic, while the other remains connected
and gains a cycle. The new boundary has size three, a contradiction.
A vertex with three cut edges would make its connected shore a singleton,
contrary to cyclicity.
Thus $X,Y$ are normalized four-poles.
Their positive ranks $h,k$ satisfy $r=h+k+3$.

If $h,k\ge H_0$, apply Theorem~\ref{thm:four-cut-large-shores}.
Otherwise reorder so that $1\le h<H_0$ and $h\le k$.
If $k$ is bounded by any fixed threshold needed below, all ranks are bounded.
Increase $K_4$ to cover those cases by the trivial cycle bound.
Assume henceforth that $k$ exceeds all these absolute thresholds.

Suppose some terminal pair of $Y$ has at least $0.51\,2^k$ paths.
Suppress the other two ports and apply
Theorem~\ref{thm:high-path-cycle-tradeoff} to obtain
$c(Y)\le K_{1/100}2^{k/2}$.
Any loop created by suppression would force path density at most one half,
so it cannot occur in this case.
Every path count of $X$ is at most $2^h$.
Every forest coefficient $\sum_\alpha f_\alpha-f_\beta$ is also at most $2^h$.
The exact gluing formula and the aggregate theorem therefore give
\[
 c(G)\le c(X)+c(Y)+2^h
 \left(4\,2^k+A_4(k+3)^4 2^{k/2}\right).
\]
The first crossing term is $2^{r-1}$.
Since $h<H_0$, all remaining terms are
$O(r^4 2^{r/2})$, with an absolute implied constant.

It remains to suppose that every path density of $Y$ is at most $0.51$.
The established asymptotic cycle bound gives
\[
 c(Y)\le\frac9{16}2^k
\]
for sufficiently large $k$.
This also applies to multigraphs after uniform subdivision.
If $h\ge2$, the first bound in Lemma~\ref{lem:four-pole-low-state}
limits crossing cycles to $3.82\,2^{h+k}$.
Their deficit from $2^{r-1}$ is at least
\[
 0.18\,2^{h+k}\ge0.72\,2^k.
\]
It absorbs $c(Y)\le(9/16)2^k$ and $c(X)<2^h$ for large $k$,
because $h<H_0$ is bounded.
Thus $c(G)<2^{r-1}$ in this case.

Finally suppose $h=1$.
The degree sum gives four vertices in $X$, all of degree two.
Since $X$ is connected and loopless, it is a four-cycle.
Its six path counts are all two.
Its normalized forest vector has entries $1/2,1/2,0$.
Let $\gamma$ be the pairing with zero entry.
The exact normalized crossing contribution is
\[
 S+\frac12F+\frac12u_\gamma,
\]
where the state belongs to $Y$.
Lemma~\ref{lem:four-pole-low-state} bounds this expression by $3.7$.
The crossing deficit is therefore at least $0.6\,2^k$.
This absorbs $c(Y)\le(9/16)2^k$ and $c(X)=1$ for sufficiently large $k$.
The theorem follows after increasing $K_4$ over the bounded cases.
\end{proof}

\paragraph{Scope.}
Cyclicity is essential to the cut hypothesis.
An acyclic two-vertex shore has a four-edge boundary in every cubic graph.
The theorem does not prove an unrestricted half-exponent bound.
After the earlier reductions, the remaining cubic kernels are cyclically
five-edge-connected.
